% !TEX program = xelatex
\documentclass[11pt,a4paper]{article}

\usepackage[a4paper,margin=22mm]{geometry}
\usepackage{xeCJK}
\usepackage{fontspec}
\usepackage{amsmath,amssymb}
\usepackage{booktabs}
\usepackage{tabularx}
\usepackage{array}
\usepackage{enumitem}
\usepackage{xcolor}
\usepackage{hyperref}

\setCJKmainfont{Hiragino Mincho ProN}
\setCJKsansfont{Hiragino Sans}
\setCJKmonofont{Hiragino Sans}
\setmonofont{Menlo}[Scale=0.82]

% Route circled numbers (①..⑭), dashes, quotes, ellipsis to the CJK font
% (Hiragino has them; the Latin roman font does not).
\xeCJKDeclareCharClass{CJK}{"2460 -> "24FF, "2013 -> "2015, "2018 -> "201F, "2026, "2500}

\hypersetup{colorlinks=true, linkcolor=black, urlcolor=blue!60!black,
  pdftitle={Xinu メッシュを AIPL 駆動リアルタイム OS 化する計画と課題},
  pdfauthor={児玉靖司}}

\newcommand{\code}[1]{\texttt{#1}}
\newcommand{\prob}[1]{\textbf{\textcolor{red!55!black}{#1}}}

\title{\textbf{Xinu メッシュを AIPL 駆動リアルタイム OS 化する計画と課題}\\[2mm]
  \large フィジカル AI（アーム型ロボット）向け分散 RTOS への道筋\\[1mm]
  \normalsize ―― 3 世代 Raspberry Pi クラスタの実測を踏まえた検討 ――}
\author{児玉靖司\\\small 法政大学経営学部}
\date{2026 年 7 月 20 日}

\begin{document}
\maketitle

\begin{abstract}
\noindent
ベアメタル Embedded Xinu を、AIPL（自作アクター言語）で記述する\textbf{フィジカル AI 向け
リアルタイム OS（RTOS）}へ発展させ、最終的にアーム型ロボットを\textbf{12 台の Xinu ノード}で
制御する構想がある。本書は、その前段として\textbf{現状の Xinu を RTOS 化するうえでの改良余地}を、
本セッションで得た 3 世代 Raspberry Pi（A53/A72/A76）クラスタの実測に基づいて洗い出し、
問題点・原因・対策・優先順位を整理する。結論として、現行ポートは\textbf{協調型スケジューラ}・
\textbf{不安定なタイマ割り込み}・\textbf{GC やデマンドページングによる非有界遅延}という、RTOS の
前提を欠く箇所を持つ。一方 Xinu の\textbf{小ささ}は WCET 解析・安全性の見通しにおいて強みであり、
「巨大 RTOS を削る」より「小さい Xinu に実時間保証と AIPL 型を足す」方が、AIPL 一貫路線に整合し、
差別化になる。本書は 4 段階のロードマップと、各段階の受入基準（実測ジッタ等）を示す。
\end{abstract}

\tableofcontents
\newpage

% ======================================================================
\section{背景と目標}

\subsection{目標像}
\begin{itemize}[leftmargin=1.4em,noitemsep]
  \item \textbf{言語}：制御・知覚・計画を原則 AIPL（型推論・効果推論を持つアクター言語）で記述する。
  \item \textbf{OS}：Xinu を土台に、\textbf{ハード実時間}（周期制御ループ、境界化された割り込み遅延）を
        提供する RTOS 化する。
  \item \textbf{対象}：アーム型ロボット等のフィジカル AI。最終形は\textbf{12 ノードの Xinu} を
        接続し、関節・センサ・計画をノードに割り当てる分散制御。
\end{itemize}

\subsection{本書の立場}
「12 台接続」に進む前に、\textbf{1 ノードの Xinu が RTOS 要件を満たすか}を先に固める。分散化は
その上に載る。以下では単一ノードの RTOS 課題（第 \ref{sec:rtos} 節）と、分散・並列計算機としての
課題（第 \ref{sec:dist} 節）を分けて述べる。

% ======================================================================
\section{現状（本セッションの実測から）}
\label{sec:now}

3 世代 Pi（A53=Pi3 / A72=Pi4 / A76=Pi5）で、ワーカープール型 SMP・キャッシュ実験・WiFi メッシュを
実機検証した。RTOS 化の観点で関係する事実を挙げる。

\begin{table}[htbp]\centering
\caption{RTOS 化に関係する現状の実測事実}
\small
\begin{tabularx}{\textwidth}{@{}lX@{}}
\toprule
観測 & 内容と含意 \\
\midrule
スケジューラ & 起動ログ \code{sched: cooperative ctxsw}。\prob{協調型（非プリエンプティブ）}。 \\
タイマ割り込み & \code{clktime} が常に 0 を返す＝\prob{周期割り込みが安定発火せず}。フリーラン
  カウンタ直読で回避したが、RTOS には割り込み遅延・ジッタの保証が要る。 \\
マルチコア & ワーカープール型：core0 が全 OS＋制御、core1--3 は IRQ マスクの計算専用。
  runtime は\prob{非リエントラント}＝対称 SMP 不可。計算律速は \(\times 3.99\) スケール。 \\
D-cache & 既定 OFF（コヒーレンシのため）。本セッションで\textbf{明示キャッシュ管理}により多コア
  D-cache を 3 世代で成立させ、単一コア再帰処理が \(\times 14\)--\(\times 16\) 高速化。 \\
メモリ & AIPL に\prob{約 8 秒周期のアクタ GC}、Pi5 に\prob{デマンドページング}——いずれも
  ハード実時間の非有界遅延源。 \\
待ち機構 & ワーカープールの \code{SMP\_WAIT\_LIMIT}=\(2\times10^9\) スピン等、\prob{非有界の待ち}が散在。 \\
ネットワーク & WiFi IBSS メッシュ（\code{MANET} 10.0.0.$n$）＋MANET/UDP。本セッションでは各台
  参加後も \prob{近隣数 0（セル未収束）}＝WiFi メッシュは\textbf{非決定的}で制御平面に不適。 \\
\bottomrule
\end{tabularx}
\end{table}

% ======================================================================
\section{RTOS 化の課題（単一ノード）}
\label{sec:rtos}

\subsection{致命的ギャップ：スケジューリングとタイミング}

\begin{enumerate}[leftmargin=1.7em]
  \item[\textbf{①}] \prob{協調型スケジューラ}。ロボット制御（例：1\,kHz サーボループ）は
    \textbf{固定優先度プリエンプティブ}（Rate\,Monotonic、必要なら EDF）が必須。共有資源
    （ロックフリー・メールボックス等）で\textbf{優先度逆転}が起こるため\textbf{優先度継承}を要する。
  \item[\textbf{②}] \prob{タイマ割り込みの不安定}。\code{clktime=0} は周期割り込みが安定発火して
    いなかった証拠。RTOS では周期 tick と\textbf{割り込み遅延（latency＋jitter）を実測・上限保証}
    しなければ成立しない。精密デッドラインには tickless＋one-shot タイマ。
\end{enumerate}

\subsection{決定性を壊す要素の除去}

\begin{enumerate}[leftmargin=1.7em]
  \item[\textbf{③}] \prob{GC の停止}。AIPL の周期 GC は制御ループを殺す。RT パスは\textbf{無割当}
    （アクタ・メッセージを事前確保）、または\textbf{インクリメンタル／リージョン GC}。GC する
    best-effort 領域と GC しない RT 領域を分離する。
  \item[\textbf{④}] \prob{デマンドページング}。ページフォルト＝非有界遅延。RT タスクのメモリは
    \textbf{ピン留め（フォルト禁止）}にする。
  \item[\textbf{⑤}] \prob{非有界のスピン／タイムアウト}。\code{SMP\_WAIT\_LIMIT} 等の巨大ループは
    ジッタ源。すべての待ちを\textbf{有界・イベント駆動}へ。
\end{enumerate}

\subsection{実時間 I/O とマルチコアの使い方}

\begin{enumerate}[leftmargin=1.7em]
  \item[\textbf{⑥}] \textbf{RT ドライバ級の新設}。現行 USB/net は best-effort。ロボット用に
    PWM／ステップ生成（ハードタイマ）・エンコーダ捕捉・I2C/SPI/CAN の\textbf{有界トランザクション}を
    「WCET 保証つき RT ドライバ」として分離する。
  \item[\textbf{⑦}] \textbf{AMP（非対称）パーティショニング}。非リエントラント runtime ゆえ対称 SMP は
    不可。コア／ボードに\textbf{役割固定（制御コア／センサコア／計画コア）}する AMP が適合し、
    「12 Xinu」構想（関節・サブシステムごと独立 RT パーティション）と一致する。制御コアは裸で隔離。
\end{enumerate}

\subsection{安全（フィジカル AI ＝ 人・機械を傷つけ得る）}

\begin{enumerate}[leftmargin=1.7em]
  \item[\textbf{⑧}] \textbf{ウォッチドッグ・デッドラインミス→セーフ状態}。締切超過を検出したら
    モータを安全トルク／ブレーキへ。1 ノード障害を波及させない\textbf{フォールト封じ込め}。現在未整備。
  \item[\textbf{⑨}] \textbf{Capability＋MMU 隔離}。既存構想の Capability を拡張し、\textbf{MMU 領域で
    アクタを隔離}（MMU は既に稼働）。バグった AIPL アクタが制御ループを壊せないようにする
    （「そのアクタは自関節しか触れない」）。
\end{enumerate}

\subsection{AIPL を「実時間言語」にする（差別化点）}

\begin{enumerate}[leftmargin=1.7em]
  \item[\textbf{⑩}] \textbf{タイミングを型／効果で表現し、スケジューラビリティを静的検証}。アクタに
    周期・デッドライン・WCET を宣言 → コンパイラが\textbf{優先度付き RT タスクへ写像}。
    \textbf{効果型で「この RT アクタは割当／ブロック／GC／send をしない」を証明}（③④⑤を言語で担保）。
    Rate\,Monotonic 上限で\textbf{スケジュール可能性を静的判定}する。これは「AIPL で書く RTOS」を
    他にない形で強くする。
\end{enumerate}

% ======================================================================
\section{分散・並列計算機としての課題}
\label{sec:dist}

12 ノードを「全体として 1 台の並列計算機／ロボット」として使うための課題。

\begin{enumerate}[leftmargin=1.7em]
  \item[\textbf{⑪}] \prob{決定的インターコネクト＋時刻同期}。本セッションで WiFi IBSS メッシュは
    参加後も\textbf{近隣数 0（セル未収束）}で、\textbf{非決定的}。制御平面は
    \textbf{CAN／EtherCAT／TSN Ethernet}、WiFi は非 RT の副系に。多関節協調には
    \textbf{PTP（IEEE\,1588）級の時刻同期}が要る（現行の「ACK 方式で時刻同期不要」は throughput 向け）。
  \item[\textbf{⑫}] \textbf{異種対応の分散スケジューラ（動的・2 階層）}。A76$>$A72$>$A53 の実測速度で
    取り分を按分し、\textbf{クロスボード work-stealing}（空いた板が忙しい板から盗む）。オンボードの
    ワーカープール＋Pi3 の least-loaded 分配を一般化。\textbf{Mac 中心の星型（SPOF）を P2P 化}。
  \item[\textbf{⑬}] \textbf{集団通信（MPI 相当）}。\code{broadcast/scatter/gather/reduce/barrier} を
    実装。木構造の scatter/reduce により台数増加でも効く。持続接続バイナリ RPC で HTTP/1.0 の
    毎回接続オーバヘッドを排す。
  \item[\textbf{⑭}] \textbf{タスク粒度}。オンボード損益分岐は 5--11\,\textmu s だが、\textbf{ネット越えは
    ms オーダ}。クロスボードで配るタスクは\textbf{$\ge$約 100\,ms の純計算}でないと損。D-cache 本番化
    （残課題は USB DMA バウンス）で計算/通信比が上がり、より小さいタスクでも分散が黒字化する。
\end{enumerate}

% ======================================================================
\section{ロードマップ（優先順）}

\begin{table}[htbp]\centering
\caption{段階的計画と受入基準}
\small
\begin{tabularx}{\textwidth}{@{}p{3.2em}Xp{10.5em}@{}}
\toprule
段階 & 内容 & 受入基準（測って合否） \\
\midrule
\textbf{P1} & \textbf{プリエンプティブ固定優先度スケジューラ＋タイマ IRQ 堅牢化}（①②）。
  固定周期タスクの release jitter を \textmu s で計測するハーネスを併設。 &
  1\,kHz 周期タスクの release jitter が上限内（例 $<100$\,\textmu s）、割り込み遅延を実測・提示。 \\
\textbf{P2} & \textbf{決定性の除去要素を潰す}：RT パス無割当＋GC 分離、ページピン、
  非有界待ちの撲滅（③④⑤）。 & RT パスで GC/フォルト/非有界待ちが 0（効果型または計測で確認）。 \\
\textbf{P3} & \textbf{AMP パーティショニング}で制御コアを隔離（⑦）＋\textbf{安全}（WD／セーフ状態、
  Capability＋MMU 隔離）（⑧⑨）。 & 制御コアが best-effort 負荷から隔離、故意フォールトで
  セーフ状態へ遷移。 \\
\textbf{P4} & \textbf{決定的ノード間通信＋時刻同期}（⑪）、\textbf{分散スケジューラ／集団通信}
  （⑫⑬⑭）、\textbf{AIPL の timing/effect 型}（⑩）を並行整備し、12 ノードへ拡張。 &
  2 ノード協調運動のジッタ上限、reduce の台数スケール、AIPL でのスケジューラビリティ静的判定。 \\
\bottomrule
\end{tabularx}
\end{table}

\noindent
\textbf{投資対効果が最も高いのは P1}である。「プリエンプティブ化＋タイマ／割り込み遅延の実測」で
\textbf{制御ループのジッタを数値で押さえる}ことが、RTOS を名乗るための最初の一歩であり、
以降のすべての基礎になる。

% ======================================================================
\section{リスクと未解決問題}

\begin{itemize}[leftmargin=1.4em]
  \item \textbf{WiFi メッシュの非決定性}：本セッションで実証（近隣 0・未収束）。制御に使うなら
    有線 RT フィールドバスへ移す判断が要る。WiFi は監視・非 RT データに限定。
  \item \textbf{D-cache 本番化の壁}：多コア D-cache は成立させたが、本番投入には USB DMA バッファの
    出自（ヒープ／スタック由来）を 64\,B 境界化・バウンスする改修が残る。
  \item \textbf{AIPL runtime の RT 化}：GC 分離と効果型による無割当保証は言語・処理系両面の設計課題。
  \item \textbf{認証・保証}：フィジカル AI では WCET と安全の\textbf{論証}が要る。Xinu の小ささは
    有利だが、AIPL 生成コードまで含めた WCET 解析の方法論が要る。
\end{itemize}

% ======================================================================
\section{結論}

現行の Xinu ポートは、\textbf{協調型スケジューラ・不安定タイマ・GC／ページングによる非有界遅延}
という点で、まだ RTOS 要件を満たさない。だが Xinu の\textbf{小ささ}は WCET 解析・安全論証に有利で、
\textbf{「小さい Xinu に実時間保証と AIPL 型を足す」}方針は、御社の AIPL 一貫路線に整合し差別化になる。
まず\textbf{P1（プリエンプティブ化＋割り込み遅延の実測）}に着手し、制御ループのジッタを数値で
押さえることを推奨する。分散（12 ノード）化は、単一ノードの実時間保証を固めた上で、決定的
インターコネクトと動的分散スケジューラの上に載せるのが順序として正しい。

\end{document}
